\chapter{Обзор Викиданных}
\label{ch:ReviewAboutWD}

\section{Викиданные}

Викиданные~--- это структурированная и совместно редактируемая база данных\footnote[][-15pt]{Викиданные (Wikidata)~--- это свободная, совместно наполняемая, многоязычная, вторичная база данных, в которой собрана структурированная информация для поддержки работы Википедии, Викисклада и других проектов Викимедиа.}. Проект был официально запущен 30 октября 2012 года, его разработка ведётся под руководством Wikimedia Deutschland\footnote[][5pt]{Немецкое отделение Фонда Викимедиа.}. Проект создавался за счёт пожертвований Allen Institute for Artificial Intelligence, Gordon and Betty Moore Foundation и Google. Викиданные~--- это бесплатная и свободная база знаний, которая может использоваться и редактироваться людьми и компьютерными программами\autocite{Vrandecic}.\begin{marginfigure}[0.0cm]
{
	\setlength{\fboxsep}{0pt}%
	\setlength{\fboxrule}{1pt}%
	\fcolorbox{gray}{white}{\includegraphics{./review/Wikidata-logo-en.png}}
}
\caption
{Логотип Викиданных. \newline
Wikimedia Commons / \href{https://commons.wikimedia.org/wiki/File:Wikidata-logo-en.svg}{Planemad}. 
}
\label{fig:seyu}
\end{marginfigure}

Содержимое Викиданных распространяется по лицензии Creative Commons CC0, которая позволяет повторно использовать информацию самыми разными способами: пользователи могут копировать, изменять, распространять и обрабатывать эти данные в любых целях. Ещё одна особенность Викиданных~--- это многоязычность. Любой человек может редактировать Викиданные более чем на 350 языках.

Викиданные постоянно обновляются, добавляются новые объекты. На 2021 год насчитывается более 95 млн страниц и более полутора млрд правок\footnote[][5px]{Данные взяты с официальной страницы статистики Викиданных: \href{https://www.wikidata.org/wiki/Wikidata:Statistics}{https://www.wikidata.org/wiki/Wikidata:Statistics}}. В 2019 года в Викиданных было совершено более 800 тысяч правок, что превзошло количество правок в Английской Википедии и сделало Викиданные наиболее редактируемым сайтом Викимедиа\footnote[][5px]{Веб-сайт Викиданных: \href{https://www.wikidata.org/}{https://www.wikidata.org/}}.

\section{Wikidata Query Service}
\label{sect:WDQS}

Любой объект Викиданных имеет свой уникальный идентификатор и свойства. Эта информация может быть обработана с помощью компьютера, и при этом она наглядно представлена и понятна пользователям без специальной предобработки. Сайт Викиданных содержит сервис ``Wikidata Query''\footnote{Полное название инструмента ``Wikidata Query Service'', кратко~--- WDQS. Ссылка: \href{https://query.wikidata.org/}{https://query.wikidata.org/}}, включающий набор инструментов для построения SPARQL-запросов и их визуализации в виде таблиц, диаграмм, графов или географических карт.

\section{Об исследовании Викиданных}

В работе ``A large-scale collaborative ontological medical database''\autocite{Collaborative_ontological_database} 
описываются плюсы использования Викиданных для создания крупномасштабной 
совместно используемой медицинской базы данных. 
Основные требования к создаваемой базе данных таковы: 
это должна быть платформа с обновлением в реальном времени, 
с лицензией, разрешающей дальнейшее использование полученной информации, 
с возможностью редактирования на любом языке и с открытым доступом. 
Именно это и есть основные характеристики Викиданных. 
Во-первых, Викиданные~--- это открытая, редактируемая база знаний. 
Любой пользователь без навыков программирования может вносить изменения 
более чем на 350 языках и диалектах. 
Во-вторых, информация постоянно обновляется, добавляются новые объекты. 
На 2021 год Викиданные насчитывают более 18 тыс. редакторов\footnote{Для сравнения число активных редакторов в Русской Википедии 
составляет примерно 4 тыс. пользователей 
(\href{https://stats.wikimedia.org/\#/ru.wikipedia.org}{https://stats.wikimedia.org/\#/ru.wikipe-dia.org}), 
на Викискладе~--- более 15 тыс. пользователей 
(\href{https://stats.wikimedia.org/\#/commons.wikimedia.org}{https://stats.wikime-dia.org/\#/commons.wikimedia.org}).}.
В-третьих, лицензия Creative Commons CC0 обеспечивает широкое использование полученной информации. 

На самом деле у Викиданных сейчас нет конкурентов. Но принято указывать аналоги и альтернативы. Укажем и мы. Есть несколько альтернативных баз знаний:
\begin{enumerate}
\item Cyc~--- это проект компании Cycorp (Остин, США) по созданию онтологической базы знаний, позволяющий решать задачи из области искусственного интеллекта\autocite{Cyc}. Сейчас Cyc имеет исследовательскую лицензию ResearchCyc. У этой базы есть некоторые недостатки: сложность системы (сложность добавления данных
вручную), недостаток документации для изучения системы, неполнота системы.
\item Evi (ранее True Knowledge\autocite{True_Knowledge})~--- это технологическая компания в Кембридже (Англия), которая специализируется на базе знаний и программном обеспечении \textit{семантического поиска}\footnote[][5pt]{Семантический поиск \index{Информатика!Семантический поиск} — способ и технология поиска информации, основанная на использовании контекстного (смыслового) значения запрашиваемых фраз, вместо словарных значений отдельных слов или выражений при поисковом запросе.} Добавление информации в базу знаний осуществляется двумя способами: импорт из <<заслуживающих доверия>> внешних баз данных (например: Википедия) и из представления пользователей в соответствии с единообразным форматом и подробным процессом ввода. Как и в Википедии, пользователь может изменять
данные, <<соглашаться>> или <<не соглашаться>> с информацией, представленной True Knowledge. Система может отклонить любые факты, которые семантически несовместимы с другими утверждениями, в отличие от Викиданных, где могут
храниться противоречивые данные.
\item DBPedia\index{Программное обеспечение!DBPedia}~--- краудсорсинговый проект, направленный на извлечение структурированной информации из данных, созданных в рамках проекта Википедия и публикации её в виде доступных под свободной лицензией наборов данных. Проект был отмечен как один из наиболее известных примеров реализации концепции связанных данных.
Проект был начат группой добровольцев из Свободного университета Берлина и Лейпцигского университета, в сотрудничестве с фирмой OpenLink Software, первый набор данных опубликован в 2007 году. С 2012 года активным участником проекта является Университет Мангейма.
\end{enumerate}
В Викиданных информация представлена в виде объектов (или элементов), связанных между собой с помощью свойств\footnote{ Например, существуют такие свойства: экземпляр \wdProperty{31}{}, подкласс \wdProperty{279}{}, часть \wdProperty{361}{}, имеет часть \wdProperty{527}{}.}. Мощь Викиданных в её большом объёме и в удивительно бустром росте и самоорганизации; в том, что к этой <<живой>> базе знаний можно обращаться с помощью  SPARQL-запросов, представлять результаты их выполнения в виде таблиц, графов, диаграмм или сохранять в нужном формате (CSV, JSON, SVG).

Викиданные могут взять на себя роль централизованного хранилища данных. В статье\autocite{Falcon_2.0} приводится пример использования Викиданных в качестве централизованной и общедоступной базы знаний для системы FALCON 2.0. Эта система идентифицирует сущности в коротком тексте или вопросе, а затем связывает их ссылками с соответствующими объектами Викиданных.

\section{Неоднозначность объекта Викиданных}

Любой объект Викиданных имеет свойства. Одно из них~--- это <<экземпляр класса>> (\wdProperty{31}{instance of}). Оно определяет класс, к которому принадлежит объект. Мы обнаружили, что один объект Викиданных может соответствовать нескольким классам.
Некоторые объекты являются экземплярами совершенно разных классов. Например, \wdqName{Королевская шведская академия наук}{191583} является экземплярами сразу трёх классов: академии наук, сооружения и королевской академии Швеции. Такое определение классов имеет место быть, поскольку этот объект можно рассматривать и как организацию, целью которой является развитие науки, и как архитектурное сооружение. 

Этот пример относится к задаче разрешения лексической многозначности или WSD-задаче. Многие ученые занимались данной задачей, в том числе итальянская учёная Angela Fogarolli. В работе\autocite{Fogarolli} она выделяет объекты неоднозначностей, которые соответствует нескольким классам в зависимости от контекста и допускает наличие нескольких классов в свойстве ``instance of''.

\section{Качество Викиданных и её техническая платформа}

Викиданные существуют с 2012 года. На 2021 год редакторами проекта являются более 200 тыс. пользователей, которые сделали более 50 млн правок.

В диссертации Alessandro Piscopo\autocite{Piscopo} рассказывается о социально-технических процессах и качестве данных проекта Викиданные, о том, что пользователи Викиданных имеют возможность добавлять отдельные фрагменты информации, выполнять редактирование через различные интерфейсы и работать с такими платформами как Википедия, но при этом они в полной мере несут ответственность за поддержание схемы графа знаний\autocite{Knowledge_Graphs} в рабочем состоянии. Однако эту работу должна выполнять команда обученных специалистов в соответствии с чётко продуманными методами. Эти действия осуществляются с помощью инструментов, которые составляют техническую основу системы.

Особым инструментом как в Викиданных, так и в Википедии являются \textit{боты}\footnote{Более подробно о ботах см. в главе <<Боты в Викиданных>> на стр. \pageref{ch:bots}.}. Это части программного обеспечения, которые автоматически могут выполнять различные действия на платформе с большой скоростью (более тысячи правок в минуту). Их основная задача~--- это редактирование существующих данных, добавление и импорт новых даных из других ресурсов. Боты создают отчёты, с помощью которых пользователь может исправлять некоторые неточности. 

Таким образом, боты являются одним из ключевых технических компонентов Викиданных. Пользователи добавляют и модифицируют данные, а также общаются между собой на форуме и на страницах обсуждений объектов Викиданных. Также доступны плагины, которые предупреждают редакторов, когда они собираются выполнить ревизию, которая может привести к каким-либо ошибкам в данных.

В статье <<Сетевая структура научных революций>>\autocite{Network_structure_revolutions} на примере Википедии рассматривается процесс формирования знаний в виде постоянно растущих сетей из статей и связывающих их гиперссылок. Эта концепция реализуется за счет заполнения пробелов в знаниях. Авторы сформулировали цель своей работы в одном предложении: <<Авторы проверяют теории научного прогресса на растущих концептуальных сетях и раскрывают управляемые данными условия, лежащие в основе прорывов>>\footnote{Оригинальный текст (англ.):  ``The authors test theories of scientific progress on growing concept networks and reveal data-driven conditions underlying breakthroughs''.},\autocite{Network_structure_revolutions}.

В процессе исследований (Zhou, Ju и Blevins, 2020)\autocite{Network_structure_revolutions}  было проведено ранжирование всех статей Википедии в виде сети по определённым критериям. Каждый узел сети соответствует определённой статье, имя узла --- это заголовок статьи, год рождения узла~--- это первый год, указанный во введении или в разделе истории как год, когда концепция была задумана. Затем на основе текущего состояния сетей были определены некоторые закономерности в эволюции этих структур на протяжении времени и периоды, когда сеть наиболее быстро менялась.

Полученные результаты показали, что человеческие знания растут и как следствие происходит постепенное изменение сетевой структуры (заполняются некоторые пробелы в знаниях). Авторы статьи\autocite{Network_structure_revolutions} считают, что знания, обнаруженные при заполнении пробелов, будут иметь важное значения для научных инноваций. 

Эта статья имеет непосредственное отношение к качеству Викиданных, потому что информация для Викиданных чаще всего берётся из Википедии. Если будут заполнены пробелы в Википедии, то новые данные обязательно будут добавлены в Викиданные, и база знаний станет более полной и подробной.
